\section{Conclusiones}

\subsection{Sobre el proceso}

\textbf{Primera: instalar el JDK y configurarlo son dos tareas distintas, y la segunda
es la que realmente se estaba practicando.} La instalación se resuelve con un solo
comando de APT y difícilmente falla. Lo que sí falla —y lo que ninguna herramienta
avisa— es la configuración de las variables de entorno: un sistema puede tener el JDK
perfectamente instalado, responder correctamente a \texttt{java -version} y aun así ser
incapaz de compilar un proyecto con Maven, porque \texttt{JAVA\_HOME} no existe o apunta
a otro lado.

\textbf{Segunda: la distinción entre JRE y JDK explica el error más frecuente.} El
estado inicial del equipo documentado en este reporte —una JVM funcionando sin
compilador— es exactamente el que produce la confusión de creer que "Java ya está
instalado". La comprobación que separa un caso del otro no es \texttt{java -version}
sino \texttt{javac -version}.

\textbf{Tercera: \texttt{JAVA\_HOME} debe apuntar a un directorio, no a un ejecutable.}
Es el error conceptual que más se repite. La variable nombra el \textbf{directorio raíz}
de la instalación —el que contiene \texttt{bin}, \texttt{lib} y \texttt{conf}—, porque
las herramientas que la consultan construyen a partir de ella rutas como
\texttt{\$JAVA\_HOME/bin/javac} o \texttt{\$JAVA\_HOME/lib/tools.jar}. Apuntarla al
binario rompe todas esas construcciones a la vez.

\textbf{Cuarta: la persistencia depende del archivo elegido, no del comando.} Escribir
\texttt{export JAVA\_HOME=...} directamente en la terminal funciona hasta que se cierra
la ventana. Que la configuración sobreviva depende de haberla escrito en el archivo
correcto, y ese archivo cambia según quién necesite ver la variable: el usuario en su
terminal, las aplicaciones gráficas, o los servicios del sistema.

\subsection{Errores posibles durante el proceso y su solución}

El siguiente catálogo reúne los problemas que pueden presentarse, ordenados por
frecuencia, con su causa y la forma de corregirlos.

\rowcolors{2}{rowblue}{white}
\footnotesize
\renewcommand{\arraystretch}{1.15}
\setlength{\tabcolsep}{4pt}

\begin{longtable}{|>{\raggedright\arraybackslash}p{4.1cm}|>{\raggedright\arraybackslash}p{4.6cm}|>{\raggedright\arraybackslash}p{5.5cm}|}
	\caption{Errores frecuentes en la instalación y configuración del JDK}
	\label{tab:errores} \\
	\hline
	\rowcolor{headerblue}
	\headercell{Síntoma} & \headercell{Causa} & \headercell{Solución} \\
	\hline
	\endfirsthead
	\hline
	\rowcolor{headerblue}
	\headercell{Síntoma} & \headercell{Causa} & \headercell{Solución} \\
	\hline
	\endhead

	\texttt{javac: command not found} pero \texttt{java} sí funciona & Se instaló el
	paquete \texttt{jre} en lugar del \texttt{jdk} & \texttt{sudo apt install
	openjdk-21-jdk}. El sufijo del paquete es lo que distingue uno de otro. \\ \hline

	\texttt{echo \$JAVA\_HOME} sigue vacío tras editar el archivo & No se releyó la
	configuración: la terminal abierta conserva el entorno anterior & \texttt{source
	\textasciitilde/.bashrc}, o cerrar y abrir la terminal. \\ \hline

	\texttt{\$JAVA\_HOME/bin/java: No such file or directory} & La variable apunta al
	ejecutable o al subdirectorio \texttt{bin} en vez de a la raíz del JDK & Usar
	\texttt{/usr/lib/\allowbreak jvm/\allowbreak java-21-\allowbreak openjdk-amd64}, sin \texttt{/bin} ni \texttt{/bin/java}
	al final. \\ \hline

	Maven o Gradle reportan \texttt{JAVA\_HOME is not set} aunque \texttt{java} funcione
	& La variable se definió solo para terminales interactivas y el proceso no la hereda
	& Definirla en \texttt{\textasciitilde/.profile} o en \texttt{/etc/profile.d/}, y
	cerrar la sesión para que surta efecto. \\ \hline

	La configuración se pierde al cerrar la terminal & Se ejecutó \texttt{export} en la
	línea de comandos, sin escribirlo en ningún archivo & Agregar las líneas a
	\texttt{\textasciitilde/.bashrc} y volver a cargarlo. \\ \hline

	Tras editar el \texttt{PATH}, ningún comando funciona (\texttt{ls}, \texttt{cd}\dots)
	& Se sobrescribió el \texttt{PATH} en vez de ampliarlo & Recuperar la sesión con
	\texttt{export PATH=/usr/bin:/bin} y corregir la línea para que termine en
	\texttt{:\$PATH}. \\ \hline

	\texttt{java -version} reporta una versión distinta a la instalada & Hay varias
	versiones y el sistema de alternativas apunta a otra & \texttt{sudo
	update-alternatives --config java} y elegir la deseada. \\ \hline

	El IDE compila con una versión y la terminal con otra & \texttt{JAVA\_HOME} y
	\texttt{update-alternatives} apuntan a JDK distintos & Alinear ambos, y revisar
	además el JDK configurado dentro del propio IDE. \\ \hline

	\texttt{E: se interrumpió la ejecución de dpkg} en cualquier comando de APT & Una
	instalación anterior de \textbf{cualquier} paquete quedó a medio configurar, y APT se
	bloquea hasta repararlo & \texttt{sudo dpkg --configure -a}, después \texttt{sudo apt
	--fix-broken install}. Diagnosticar con \texttt{sudo dpkg --audit}. \\ \hline

	\texttt{Unable to locate package openjdk-21-jdk} & Índices de APT desactualizados, o
	la versión no existe para esa distribución & \texttt{sudo apt update} y verificar con
	\texttt{apt-cache policy}. En versiones antiguas de Ubuntu puede hacer falta un PPA.
	\\ \hline

	\texttt{404 Not Found} a mitad de la descarga & El catálogo local apunta a una
	versión que el repositorio ya reemplazó & \texttt{sudo apt update} y repetir la
	instalación. \\ \hline

	\texttt{Permission denied} al editar \texttt{/etc/environment} & El archivo es del
	sistema y requiere privilegios & Usar \texttt{sudo nano /etc/environment}. \\ \hline

	\texttt{/etc/environment} no surte efecto & Se escribió con la sintaxis
	\texttt{export CLAVE=valor} & Ese archivo solo admite pares \texttt{CLAVE="valor"},
	sin \texttt{export}. \\ \hline

	\texttt{apt autoremove} propone eliminar programas en uso & Aplicaciones de
	escritorio que dependen del Java que se está retirando & Revisar la lista antes de
	aceptar; si hay dudas, no desinstalar y convivir con ambas versiones. \\ \hline

	Todo funciona hasta reiniciar, y entonces se pierde & La variable se definió en un
	archivo que la sesión nueva no lee & Verificar el archivo elegido según la tabla
	\ref{tab:alcance} y comprobar con una terminal recién abierta. \\ \hline

	\texttt{Unsupported\allowbreak ClassVersion\allowbreak Error} al ejecutar & El \textit{bytecode} se compiló
	con una versión más nueva que la JVM que lo ejecuta & Alinear las versiones, o
	compilar con \texttt{javac --release 21}. \\ \hline
\end{longtable}

\normalsize
\setlength{\tabcolsep}{6pt}

\subsection{Recomendaciones}

\begin{enumerate}
	\item \textbf{Diagnosticar antes de instalar.} Ejecutar \texttt{javac -version} y
	\texttt{echo \$JAVA\_HOME} antes de tocar nada evita instalar sobre una configuración
	que ya existía y que quedará huérfana.

	\item \textbf{Verificar con \texttt{\$JAVA\_HOME/bin/java -version}}, no solo con
	\texttt{echo \$JAVA\_HOME}. Lo primero comprueba que la ruta \textbf{existe y es
	válida}; lo segundo solo muestra una cadena de texto que puede estar mal escrita.

	\item \textbf{Comprobar en una terminal nueva.} Es la única forma de saber si la
	configuración es persistente o si solo vive en la sesión actual.

	\item \textbf{Usar el enlace simbólico y no la ruta con número de versión.} En
	Ubuntu, \texttt{/usr/lib/\allowbreak jvm/\allowbreak java-21-\allowbreak openjdk-amd64} se mantiene estable entre
	actualizaciones menores, mientras que una ruta como
	\texttt{\dots/jdk-21.0.11} deja de existir en cuanto el paquete se actualiza a
	\texttt{21.0.12}, y \texttt{JAVA\_HOME} queda apuntando a la nada.

	\item \textbf{Para varias versiones, considerar SDKMAN.} En equipos de desarrollo que
	alternan entre proyectos con distintas versiones de Java, una herramienta de gestión
	como SDKMAN evita reconfigurar variables a mano. Conviene saber que ese tipo de
	herramientas también modifica \texttt{JAVA\_HOME}, y que mezclarla con una
	configuración manual en \texttt{.bashrc} produce conflictos difíciles de rastrear.
\end{enumerate}

\subsection{Sobre la distribución utilizada}

Se documentó la instalación de OpenJDK en lugar de Oracle JDK. Vale la pena cerrar
señalando qué habría cambiado y qué no.

\textbf{No habría cambiado nada de lo esencial.} Ambas distribuciones parten del mismo
código fuente del proyecto OpenJDK, superan el mismo conjunto de pruebas de
compatibilidad y ejecutan idénticamente el mismo \textit{bytecode}. Todo lo relativo a
\texttt{JAVA\_HOME}, al \texttt{PATH}, al sistema de alternativas y a los errores
catalogados arriba aplica exactamente igual.

\textbf{Habría cambiado el procedimiento de obtención y la ruta.} Con Oracle JDK habría
que descargar el paquete \texttt{.deb} desde el sitio de Oracle aceptando su licencia,
instalarlo con \texttt{sudo dpkg -i}, y \texttt{JAVA\_HOME} apuntaría a
\texttt{/usr/lib/jvm/jdk-21-oracle}. A cambio de ese trabajo manual se pierde la
actualización automática por APT: cada parche de seguridad exigiría repetir la descarga.

\textbf{La diferencia real es de licencia.} OpenJDK se distribuye bajo GPLv2 con
excepción de enlazado, sin restricciones de uso. Oracle JDK se rige por los
\textit{Oracle No-Fee Terms and Conditions}, que permiten el uso gratuito bajo
condiciones específicas y han cambiado varias veces en los últimos años. Para un entorno
académico y para la mayoría de los proyectos, OpenJDK es la opción con menos fricción
legal y operativa.
